\section{Discussion and open questions}
\label{sec:scope}

The query improvement rests on storing each candidate set as the
conjunction of the adjacency constraints created by earlier pivots, and on
allocating estimation error according to the cost of each level.
\Cref{tab:bounds} leaves two central questions.

\paragraph{A worst-case quantum advantage.}
Can the quantum and randomized query complexities be separated in the worst
case?  The randomized complexity on $N=2^n$ vertices lies between
$\Omega(N^{1-1/\sqrt2})$ and $O(N)$, while the quantum upper bound is
$\widetilde O(\sqrt N)$.  A randomized lower bound of
$N^{1/2+\Omega(1)}$ would establish such a separation.  The implicit-set
recursion explains the difference between the available upper bounds:
classical rejection sampling at depth $i$ needs $2^i$ membership tests,
whereas quantum search needs $2^{i/2}$.

\paragraph{The optimal quantum exponent.}
Is $\sqrt N$ optimal in the worst case?  The current quantum bounds are
$\Omega(N^{1/12})$ and $O(\sqrt N\log N\log\log N)$.  The final levels of
the recursion search sets of constant size inside a universe of size
$\Theta(N)$, but these sets are defined by $O(\log N)$ adjacency
constraints, so the optimality of Grover search does not establish a lower
bound for this relation.

\paragraph{Scope.}
The distributional separation on $G(N,1/2)$ in \cref{cor:separation} and
the multicolour bound in \cref{cor:multicolour} hold under their stated
contracts.  \Cref{app:experiments} reports small sanity-check experiments;
the query bounds rest on the proofs in the appendices.  The paper claims
query bounds only, not gate complexity or hardware performance.
